% =====================================================
% 题型：导数几何意义求切线方程选择题 (q_func_tangent_line.tex)
% 知识板块：导数切线
% 主思维：函数思维
% 副思维：定义域优先、卷面表达
% 错因标签：切点代错、导数斜率漏写、忽视定义域 x>0
% 讲义用途：导数卷面模板
% =====================================================
% 难度：★ (基础)
% =====================================================
\newcommand{\qFuncTangentLineStem}{曲线 \(y = 5x + 8\ln x\) 在点 \((1, 5)\) 处的切线方程为 \hfill （\quad）}

\newcommand{\qFuncTangentLineOptions}{%
  \mcoptions[2]{\(y = 3x + 2\)}{\(y = 5x\)}{\(y = 8x - 3\)}{\(y = 13x - 8\)}%
}

\newcommand{\qFuncTangentLineAnswer}{D}

\newcommand{\qFuncTangentLineSolution}{%
  【解题思路】\\
  利用导数的几何意义求解切线方程。首先对曲线方程求导，求出导函数 \(y'\)；然后将切点的横坐标 \(x = 1\) 代入导函数，求得切线的斜率 \(k\)；最后利用点斜式公式写出切线方程，并化简为常用形式。\\
  【解析计算】\\
  已知曲线方程为 \(y = 5x + 8\ln x\)。\\
  求导得：\\
  \(y' = 5 + \dfrac{8}{x}\)。\\
  当 \(x = 1\) 时，切线的斜率为：\\
  \(k = y'\Big|_{x=1} = 5 + 8 = 13\)。\\
  已知切点为 \((1, 5)\)，根据直线的点斜式方程，切线方程为：\\
  \(y - 5 = 13(x - 1)\)。\\
  整理化简得：\\
  \(y = 13x - 8\)。\\
  故选 D。%
}

\newcommand{\qFuncTangentLineDiagnosis}{%
  学生偶有将求导公式记错的情况（如将 \(\ln x\) 的导数记错，或常数项未导掉）。部分学生在得到斜率后套用直线方程公式时容易出现负号或常数项计算失误。同时要建立定义域意识，即使是切线，对于 \(\ln x\) 的定义域限制 \(x>0\) 必须有初步认识。%
}

\newcommand{\qFuncTangentLineTeachingNotes}{%
  切线问题是导数题中最基本、也是高考必考的题型。应引导学生规范写出三步法：①求导数；②求斜率；③写方程。并强调一定要检查点 \((1, 5)\) 是否在曲线上（本题“在点 \((1, 5)\) 处”已表明其为切点）。%
}
